% Copy this template for each proposal and replace the title, names, and prompts.
% Edit the title and author below, and their matching PDF metadata in hypersetup.
% The headings are a starting point; adapt them if another structure communicates
% your idea more clearly. This is an exploration, so say what you are still unsure about.
% Include related work where it helps explain your idea.
%
% Build from the folder containing your copy:
%   pdflatex project_outline.tex
%
% The layout and content match the Markdown template converted with:
%   pandoc project_outline.md --from=gfm+raw_attribute+smart --standalone -o project_outline.tex
% After conversion, use $...$ for inline math and equation* for unnumbered displays.
%
% Options for packages loaded elsewhere
\PassOptionsToPackage{unicode}{hyperref}
\PassOptionsToPackage{hyphens}{url}
\documentclass[
  11pt,
  letterpaper,
]{article}
\usepackage{xcolor}
\usepackage[margin=0.85in]{geometry}
\usepackage{amsmath,amssymb}
\setcounter{secnumdepth}{-\maxdimen} % remove section numbering
\usepackage{iftex}
\ifPDFTeX
  \usepackage[T1]{fontenc}
  \usepackage[utf8]{inputenc}
  \usepackage{textcomp} % provide euro and other symbols
\else % if luatex or xetex
  \usepackage{unicode-math} % this also loads fontspec
  \defaultfontfeatures{Scale=MatchLowercase}
  \defaultfontfeatures[\rmfamily]{Ligatures=TeX,Scale=1}
\fi
\usepackage{lmodern}
\ifPDFTeX\else
  % xetex/luatex font selection
\fi
% Use upquote if available, for straight quotes in verbatim environments
\IfFileExists{upquote.sty}{\usepackage{upquote}}{}
\IfFileExists{microtype.sty}{% use microtype if available
  \usepackage[]{microtype}
  \UseMicrotypeSet[protrusion]{basicmath} % disable protrusion for tt fonts
}{}
\makeatletter
\@ifundefined{KOMAClassName}{% if non-KOMA class
  \IfFileExists{parskip.sty}{%
    \usepackage{parskip}
  }{% else
    \setlength{\parindent}{0pt}
    \setlength{\parskip}{6pt plus 2pt minus 1pt}}
}{% if KOMA class
  \KOMAoptions{parskip=half}}
\makeatother
\pagestyle{empty}
\setlength{\emergencystretch}{3em} % prevent overfull lines
\providecommand{\tightlist}{%
  \setlength{\itemsep}{0pt}\setlength{\parskip}{0pt}}
\usepackage{graphicx}
\usepackage{titlesec}
\titleformat{\section}{\normalfont\large\bfseries}{}{0pt}{}
\titlespacing*{\section}{0pt}{1em}{0.35em}
\setlength{\parindent}{0pt}
\setlength{\parskip}{0.5em}
\raggedbottom
\makeatletter
\renewcommand{\maketitle}{%
  \begingroup\centering
  {\Large\bfseries\@title\par}
  \vspace{0.35em}
  {\normalsize\@author\par}
  \endgroup\vspace{0.6em}}
\makeatother
% Start each example on a new PDF page without raw LaTeX in the Markdown body.
\newif\ifoutlinefirstsection
\outlinefirstsectiontrue
\let\outlinesection\section
\renewcommand{\section}{%
  \ifoutlinefirstsection\outlinefirstsectionfalse\else\clearpage\fi
  \outlinesection}
\usepackage{bookmark}
\IfFileExists{xurl.sty}{\usepackage{xurl}}{} % add URL line breaks if available
\urlstyle{same}
% fallback for those not using the hyperref driver hyperxmp:
\makeatletter
\@ifundefined{xmpquote}{\newcommand{\xmpquote}[1]{#1}}{}
\makeatother
\hypersetup{
  pdftitle={{[}Project title{]}},
  pdfauthor={{[}Your name and netID, or all teammates' names and netIDs{]}},
  hidelinks,
  pdfcreator={LaTeX via pandoc}}

\title{{[}Project title{]}}
\author{{[}Your name and netID, or all teammates' names and netIDs{]}}
\date{}

\begin{document}
\maketitle

\section{\texorpdfstring{\emph{{[}MAE4110 Performance-driven
example{]}}}{{[}MAE4110 Performance-driven example{]}}}\label{mae4110-performance-driven-example}

\subsection{What do you want to
achieve?}\label{what-do-you-want-to-achieve}

Get the robot to backflip on command. A minimum test will be successful
sim-to-real transfer of a policy that can perform a backflip (full 360
rotation) from standing. Stretch goals will be to perform back-flips
while moving, or to do other acrobatics (e.g. barrel-rolls, etc.).

\subsection{How will you show it?}\label{how-will-you-show-it}

To assess performance, we will look at how ``clean'' the landing is:

\begin{itemize}
\tightlist
\item
  how quickly rotational momentum is dissipated and if it has
  oscillations
\item
  time from touch-down to stable standing
\item
  as well as number of foot-repositioning (lifting the foot after it has
  touched-down).
\item
  difference of the above between simulation and hardware experiments
\end{itemize}

\section{\texorpdfstring{\emph{{[}MAE5110 Contribution-driven
example{]}}}{{[}MAE5110 Contribution-driven example{]}}}\label{mae5110-contribution-driven-example}

\subsection{Context}\label{context}

Current state-of-the-art RL control for legged robots learn policies
that output the joint position targets for a PD controller, implemented
as part of the environment, of the form

\begin{equation*}
\tau = K_p (q_{\text{target}} - q) + K_d (\dot{q}_{\text{target}} - \dot{q}),
\end{equation*}

where $\tau$ is the joint torque, $K_p$ and $K_d$ are feedback
gains, and $q$ and $\dot{q}$ are the joint positions and velocities,
respectively. However, the policies typically only output the joint
position targets, and leave the velocity target
$\dot{q}_{\text{target}} = 0$. This is often justified by the
ambiguity introduced by an action space including both terms, since a
desired torque $\tau$ can be achieved by different combinations of
$q_{\text{target}}$ and $\dot{q}_{\text{target}}$.

\subsection{Contribution}\label{contribution}

We show that the an action space including both $q_{\text{target}}$
and $\dot{q}_{\text{target}}$ is ambiguous under typical reward
shaping, but that eliminating this ambiguity results in policies that
violate kinematic constraints less often, and are better able to deal
with behaviors that require both high velocity and abrupt changes in
velocity, such as coming to a quick stop from running.

\subsection{Experiment}\label{experiment}

We will establish the presence of ambiguity by comparing three cases:

\begin{enumerate}
\def\labelenumi{\arabic{enumi}.}
\tightlist
\item
  action-space $\mathcal{A}: [q_{\text{target}}]$ (baseline)
\item
  action-space
  $\mathcal{A}: [q_{\text{target}}, \dot{q}_{\text{target}}]$.
\item
  action-space $\mathcal{A}: [q_{\text{target}}]$ with
  $\dot{q}_{\text{target}}$ computed from a history of
  $q_{\text{target}}$.
\end{enumerate}

For each, we will compare the variance in policy output to variance in
torque output, and in overall reward-output.

We will then benchmark custom rewards that separately penalize
$q_{\text{target}}$ and $\dot{q}_{\text{target}}$ to disambiguate
the action space, and validate them in terms of variance ratio (see
above), and also in terms of frequency of the policy outputting
kinematically infeasible commands during walking, running, and
transitions.

% Optional: put supporting figures on a second page. To force a page break and
% include a figure, uncomment the following lines and replace the path and caption:
%
% \newpage
% \begin{figure}[htbp]
% \centering
% \includegraphics[width=0.8\linewidth]{figure.png}
% \caption{Caption explaining how the figure helps communicate your idea.}
% \end{figure}

\end{document}
